% =========================
% English Abstract
% =========================
\newpage
\chapter*{Abstract}
\addcontentsline{toc}{chapter}{Abstract}

\begin{table}[h]
    \begin{tabular}{l l}
        Thesis Title   & Simulation of Maternal Instinct \\
        Thesis Credits & 8 \\
        Candidate      & Mr.\ Napat Aeimwiratchai \\
        Thesis Advisor & Mr.\ Bawornsak Sakulkueakulsuk \\
        Program        & Bachelor of Engineering \\
        Field of Study & Robotics and Automation \\
        Department     & Robotics and Automation \\
        Faculty        & Institute of Field Robotics \\
        Academic Year  & 2025 \\
    \end{tabular}
\end{table}

Maternal care is a foundational form of altruism, yet the mechanisms by which prosocial strategies emerge and persist in artificial systems remain unclear.
This thesis presents a computational framework that couples a neuroendocrine-inspired controller with evolutionary agent-based modeling to study how caregiving-like behavior can emerge in a minimal dyadic setting.
Building on evidence for oxytocin- and cortisol-like modulation of social bonding and stress responses, we model a mother agent whose perception\,\allowbreak$\to$\,hormone\,\allowbreak$\to$\,motivation cascade produces behavior through weighted motivational competition (Forage, Care, Self, Protect).
The child agent is passive: it exposes internal need variables (hunger, warmth, injury) that the mother observes and acts on, creating an asymmetric caregiving dynamic without explicit cooperative rewards.

We run spatially explicit simulations in a $15\times15$ grid-world with periodic food spawning (food spawn interval $\in\{5,15,25\}$ ticks; shorter interval = more abundant food) and optional threats.
Motivation-weight genomes are optimized using a single-lineage $(1{+}1)$ evolutionary search evaluated over 10 stochastic rollouts per candidate, across 3000 generations.
We compare three plasticity regimes---E2 (genes only), P1 (adaptive-rate outcome plasticity), and P2 (fixed-rate outcome plasticity)---and three motivation-weight initialization priors: \texttt{anti\_maternal}, \texttt{random\_uniform}, and \texttt{pro\_maternal}.

Four experiments are reported:
(1)~characterizing maternal behavior dynamics in the training ecology;
(2)~showing that plasticity accelerates evolutionary fitness improvement (Baldwin-compatible pattern);
(3)~demonstrating that initialization priors shape weight displacement and residual failure modes after 3000 generations; and
(4)~evaluating generalization to nine unseen world$\times$food cells, showing graded transfer: survival is maintained under abundant food and low threat, but child hunger becomes the dominant failure mode under scarce food (interval~25) and high threat.

The framework bridges descriptive biology and abstract cooperation models, offering mechanistic insight into how minimal neuroendocrine architectures generate adaptive prosocial strategies.\\

\noindent \textbf{Keywords:} maternal behavior, neuroendocrine modulation, evolutionary computation, agent-based modeling, motivation weights, Baldwin effect, bio-inspired AI
